\section*{Resumen}
El presente Informe Interpretativo de Integración aborda una problemática estructural crítica en el campo del noreste de México: el aprovechamiento ilegal y sobreconcesión de recursos hídricos subterráneos mediante pozos clandestinos, articulado con el uso indiscriminado de plaguicidas de alta peligrosidad no autorizados por COFEPRIS en huertas agrícolas de Nuevo León. Se analiza el fenómeno bajo el prisma de la ética cívica contemporánea de Adela Cortina, distinguiendo la vulneración de la ética de mínimos de justicia universal frente a intereses corporativos particulares. Asimismo, se desglosa el marco normativo mexicano vulnerado (Artículos 4.° y 27 de la Constitución Política de los Estados Unidos Mexicanos, Ley de Aguas Nacionales, LGEEPA y Código Penal Federal). Con base en la teoría de los círculos de control e influencia, se diseña un Plan Estratégico de Intervención Profesional para el Ingeniero Agrónomo que comprende tres ejes: auditorías técnicas voluntarias de caudal y calibración de equipos, programas de transición hacia bioinsumos certificados, y mecanismos de trazabilidad digital y denuncia ciudadana protegida. Se concluye que la cultura de la legalidad no es una imposición coercitiva externa, sino la condición indispensable para la paz social, la inocuidad alimentaria y la supervivencia de los agroecosistemas.

\vspace{0.4cm}
\noindent\textbf{Palabras clave:} Ética cívica, cultura de la legalidad, Estado de derecho, gobernanza hídrica, plaguicidas, ingeniería agronómica, UANL.

\section{Introducción y Diagnóstico de la Problemática Agrícola}
La agricultura comercial en las cuencas de los ríos San Juan y Bravo enfrenta una crisis de sostenibilidad marcada por la escasez hídrica y la presión competitiva en los mercados agroalimentarios internacionales \parencite{cortina2013, cndh2020}. En este contexto, proliferan dos prácticas ilícitas recurrentes que socavan el Estado de derecho:
\begin{enumerate}
    \item \textbf{Extracción clandestina de agua subterránea:} Perforación de pozos agrícolas sin concesión oficial de la Comisión Nacional del Agua (CONAGUA) o alteración deliberada de los macromedidores de gasto volumétrico, agotando los acuíferos sobreexplotados de la región citrícola y norte del estado \parencite{cpeum2024}.
    \item \textbf{Uso y tráfico de plaguicidas prohibidos:} Aplicación clandestina de agroquímicos organoclorados y organofosforados de etiqueta roja prohibidos en México y en convenios internacionales (e.g. Convenio de Róterdam y Estocolmo), así como la quema o abandono de envases vacíos sin triple lavado en cauces de arroyos, violando la NOM-003-STPS-1999.
\end{enumerate}

Estas conductas generan una doble afectación: la degradación irreversible de los bienes comunes naturales y la vulneración flagrante de los derechos humanos a la salud, a un medio ambiente sano y al agua potable de las comunidades rurales aledañas \parencite{giusti2018}.

\section{Fundamentación Teórica y Ética Cívica}
\subsection{Ética de Mínimos de Justicia vs. Ética de Máximos}
Siguiendo la propuesta filosófica de Adela Cortina \parencite{cortina2013}, en una sociedad pluralista y democrática coexisten múltiples concepciones del bien y de la autorrealización personal (éticas de máximos). Sin embargo, la convivencia armónica sólo es posible si se garantiza una base común de principios universales innegociables: la \textbf{ética de mínimos o ética de la justicia}.

La extracción clandestina de agua constituye una violación directa de esta ética de mínimos, puesto que privatiza ilícitamente un bien público indispensable para la subsistencia de terceros. El productor o empresa que perfora al margen de la ley actúa bajo un esquema de egoísmo moral ("tragedia de los comunes"), maximizando su beneficio económico a corto plazo a expensas de la vulneración del derecho humano al agua de las generaciones presentes y futuras \parencite{giusti2018}.

\subsection{Cultura de la Legalidad y Deontología Agronómica}
La cultura de la legalidad se define como el imperio de la norma jurídica sustentado en la convicción ética compartida de su legitimidad y utilidad social \parencite{cndh2020}. El Código de Ética y Conducta de la Universidad Autónoma de Nuevo León \parencite{uanl_codigo2021} establece que sus egresados deben regirse por la honestidad, la responsabilidad social y el apego irrestricto a la verdad científica. Cuando un Ingeniero Agrónomo avala dictámenes técnicos fraudulentos, certifica volúmenes hídricos ficticios o prescribe agroquímicos proscritos, comete una falta deontológica grave que deslegitima a la profesión entera ante la sociedad.

\section{Marco Jurídico e Institucional Mexicano}
El abordaje del problema se sustenta en el orden constitucional y legal vigente en México:
\begin{table}[h!]
    \centering
    \small
    \begin{tabularx}{\textwidth}{lYY}
        \toprule
        \textbf{Ordenamiento Legal} & \textbf{Precepto Vulnerado} & \textbf{Sanción y Alcance Jurídico} \\
        \midrule
        \textbf{CPEUM} & Art. 4.° (Derecho al medio ambiente y al agua) y Art. 27 (Propiedad originaria de la Nación). & Nulidad de actos y responsabilidad patrimonial del Estado \parencite{cpeum2024}. \\
        \textbf{Ley de Aguas Nacionales} & Art. 119 (Explotación sin concesión o rebasando el volumen asignado). & Clausura de pozos, revocación de títulos y multas económicas de hasta 50,000 UMAS. \\
        \textbf{LGEEPA} & Art. 109 bis y 153 (Manejo integral de residuos peligrosos y sustancias tóxicas). & Clausura temporal/definitiva de predios y remediación obligatoria del sitio contaminado. \\
        \textbf{Código Penal Federal} & Art. 416 y 420 Quáter (Delitos contra el ambiente por sustancias tóxicas e invasión hídrica). & Penas corporales de 1 a 9 años de prisión para responsables materiales e intelectuales. \\
        \bottomrule
    \end{tabularx}
    \caption{Matriz del marco normativo federal aplicable a infracciones hídricas y ambientales.}
\end{table}

\section{Plan Estratégico de Intervención Profesional}
Bajo el modelo de Stephen Covey, el profesionista debe concentrar su energía en su \textbf{círculo de control e influencia directo}, evitando la inacción ante la magnitud del problema sistémico:
\begin{enumerate}
    \item \textbf{Eje 1: Auditoría Técnica y Transición Hídrica (Corto Plazo):}
    Implementación en ranchos y cooperativas de balances hídricos continuos mediante sensores telemétricos certificados y calibración estricta de boquillas y emisores de riego, logrando un ahorro de caudal del $30\%$ que desincentiva la perforación ilegal.
    \item \textbf{Eje 2: Sustitución de Agroquímicos y Gestión de Envases (Mediano Plazo):}
    Establecimiento de protocolos de Manejo Integrado de Plagas (MIP), incorporando agentes de biocontrol (microorganismos entomopatógenos como \textit{Bacillus thuringiensis}) y vinculación con el programa Campo Limpio de Amocali A.C. para el triple lavado, perforado y acopio del $100\%$ de envases vacíos.
    \item \textbf{Eje 3: Ética Corporativa y Denuncia Ciudadana Digital (Largo Plazo):}
    Capacitación deontológica a cuadrillas de campo y diseño de plataformas digitales anónimas seguras para alertar sobre irregularidades operativas ante PROFEPA y CONAGUA sin exponer la seguridad física de los técnicos.
\end{enumerate}

\section{Conclusiones}
1. La legalidad en el campo agrícola no es un obstáculo burocrático, sino el blindaje indispensable para acceder a mercados de exportación de alto valor y proteger la salud pública regional.
2. La formación humanista y técnica de la Facultad de Agronomía de la UANL capacita al egresado para erigirse en agente activo de cambio, liderando con integridad y rigor normativo.
3. El compromiso ético individual transforma el entorno productivo hacia una auténtica agricultura regenerativa sustentada en la justicia y la dignidad humana.

